% dimensional_algebra.tex
% Dimensional Algebra: From the Golden Axiom to the Quintic Barrier
% nos3bl33d / (DFG) DeadFoxGroup
% April 2026
%
% Compile: pdflatex dimensional_algebra.tex (twice for references)

\documentclass[11pt,a4paper]{article}

% -- Packages ----------------------------------------------------------------
\usepackage[utf8]{inputenc}
\usepackage[T1]{fontenc}
\usepackage{amsmath,amssymb,amsthm,mathtools}
\usepackage{geometry}
\usepackage{hyperref}
\usepackage{booktabs}
\usepackage{enumitem}
\usepackage{microtype}
\usepackage{xcolor}
\usepackage{array}
\usepackage{tikz}
\usetikzlibrary{arrows.meta}

\geometry{margin=1in}

\hypersetup{
  colorlinks=true,
  linkcolor=blue!70!black,
  citecolor=blue!70!black,
  urlcolor=blue!70!black,
}

% -- Theorem environments ----------------------------------------------------
\theoremstyle{plain}
\newtheorem{theorem}{Theorem}[section]
\newtheorem{lemma}[theorem]{Lemma}
\newtheorem{proposition}[theorem]{Proposition}
\newtheorem{corollary}[theorem]{Corollary}
\newtheorem{conjecture}[theorem]{Conjecture}

\theoremstyle{definition}
\newtheorem{definition}[theorem]{Definition}
\newtheorem{example}[theorem]{Example}
\newtheorem{remark}[theorem]{Remark}

% -- Shortcuts ---------------------------------------------------------------
\newcommand{\R}{\mathbb{R}}
\newcommand{\Z}{\mathbb{Z}}
\newcommand{\Q}{\mathbb{Q}}
\newcommand{\N}{\mathbb{N}}
\newcommand{\F}{\mathbb{F}}
\newcommand{\vp}{\varphi}
\newcommand{\ps}{\psi}
\newcommand{\eps}{\varepsilon}

\DeclareMathOperator{\Gal}{Gal}
\DeclareMathOperator{\Aut}{Aut}
\DeclareMathOperator{\SO}{SO}
\DeclareMathOperator{\SU}{SU}

% -- Status markers ----------------------------------------------------------
\newcommand{\proven}{\textsc{[P]}}
\newcommand{\derived}{\textsc{[D]}}
\newcommand{\conjectural}{\textsc{[C]}}
\newcommand{\metaphor}{\textsc{[M]}}

% =============================================================================
\begin{document}

\title{%
  \textbf{Dimensional Algebra:}\\[4pt]
  \textbf{From the Golden Axiom to the Quintic Barrier}%
}
\author{%
  nos3bl33d\thanks{%
    DeadFoxGroup (DFG). Correspondence: \texttt{nos3bl33d@dfg}.%
  }%
}
\date{April 2026}

\maketitle

% -- Abstract ----------------------------------------------------------------
\begin{abstract}
We subject five claims from the dimensional algebra framework to formal
analysis, proving two as theorems, identifying one as a precise metaphor
wrapping a real construction, and classifying two as framework conventions
requiring further axiomatization. The proven results are: (1)~the
alternating group $A_5 \cong \mathrm{Rot}(\text{Icosahedron})$ is the
smallest non-solvable group, which is why the general quintic admits no
radical solution (Abel--Ruffini--Galois--Klein, 1824--1884); and (2)~the
golden ratio $\vp = (1+\sqrt{5})/2$ is the infimum of exponential growth
rates among second-order linear recurrences with positive integer
coefficients (proved here as Theorem~\ref{thm:phi-minimal}). The
``square $+1 =$ cube'' construction is formalized as the standard
hypercube product $C_n = C_{n-1} \times [0,1]$
(Proposition~\ref{prop:hypercube}). The degree/zero correspondences
are classified as framework conventions awaiting rigorous grounding.
Each claim is tagged \proven\ (proven), \derived\ (derived under stated
assumptions), \conjectural\ (conjectural), or \metaphor\ (metaphor
encoding a real structure).
\end{abstract}

\tableofcontents

% =============================================================================
\section{Notation and Setup}\label{sec:notation}
% =============================================================================

Throughout, we fix the following constants from the dodecahedral framework:

\begin{definition}[Framework Constants]\label{def:constants}
Let $d = 3$ (vertex degree of the dodecahedron), $p = 5$ (face degree),
$\chi = 2$ (Euler characteristic of the $2$-sphere). The dodecahedron has
$V = 20$ vertices, $E = 30$ edges, $F = 12$ pentagonal faces, satisfying
$V - E + F = \chi = 2$.
\end{definition}

\begin{definition}[The Axiom]\label{def:axiom}
The \emph{golden axiom} is the equation
\begin{equation}\label{eq:axiom}
  x^2 = x + 1,
\end{equation}
equivalently $x^2 - x - 1 = 0$, with roots
\begin{equation}\label{eq:roots}
  \vp = \frac{1+\sqrt{5}}{2} = 1.618\,033\,988\,7\ldots, \qquad
  \ps = \frac{1-\sqrt{5}}{2} = -0.618\,033\,988\,7\ldots
\end{equation}
\end{definition}

\begin{remark}[Vieta Relations] \proven\
The roots satisfy $\vp + \ps = 1$ and $\vp \cdot \ps = -1$. Both are
immediate from Vieta's formulas applied to $x^2 - x - 1 = 0$.
\end{remark}


% =============================================================================
\section{Claim 1: The Square Off Center Makes a Cube}\label{sec:cube}
% =============================================================================

\subsection{The Informal Claim} \metaphor

The framework states: \emph{``A square pushed off center by $+1$ creates a
cube. $360 + 90 = 450$ degrees total.''} We formalize what is geometrically
true and identify what is metaphor.

\subsection{The Hypercube Product}

\begin{definition}[Unit Hypercube]\label{def:hypercube}
The $n$-dimensional unit hypercube is
\[
  C_n = [0,1]^n = \{(x_1, \ldots, x_n) \in \R^n : 0 \le x_i \le 1
  \ \forall\, i\}.
\]
\end{definition}

\begin{proposition}[Dimension by Cartesian Product] \proven\
\label{prop:hypercube}
For all $n \ge 1$,
\[
  C_n = C_{n-1} \times [0,1],
\]
where the product is taken along the $n$-th coordinate axis, which is
perpendicular to the hyperplane containing $C_{n-1}$. The extrusion
distance equals $1$, the side length of $C_{n-1}$.
\end{proposition}

\begin{proof}
By definition, $C_{n-1} \times [0,1] = \{(x_1,\ldots,x_{n-1},x_n) :
(x_1,\ldots,x_{n-1}) \in C_{n-1},\; x_n \in [0,1]\} = [0,1]^n = C_n$.
The $n$-th axis $\mathbf{e}_n$ is orthogonal to $\mathrm{span}
(\mathbf{e}_1,\ldots,\mathbf{e}_{n-1})$ by the standard inner product on
$\R^n$. The extrusion length is $|[0,1]| = 1 = $ side length of $C_{n-1}$.
\end{proof}

\begin{remark}[Why ``Off Center'' is Imprecise]
An extrusion of a square along a non-perpendicular direction, or for a
distance $\ell \ne 1$, produces a parallelepiped or rectangular prism, not
a cube. The construction requires:
\begin{enumerate}[nosep]
  \item The extrusion direction is \emph{perpendicular} to the existing
    $(n{-}1)$-dimensional hyperplane.
  \item The extrusion distance equals the side length of $C_{n-1}$.
\end{enumerate}
The phrase ``off center by $+1$'' is a mnemonic for these two conditions:
the $+1$ encodes the unit-length extrusion, and ``off center'' evokes
perpendicularity. As a precise geometric statement, the claim should read:
\emph{``The unit $n$-cube is the unit $(n{-}1)$-cube extruded by $1$ along
a new perpendicular axis.''}
\end{remark}

\subsection{Connection to the Axiom}

\begin{proposition}[The Axiom as a Dimensional Operator] \derived\
\label{prop:axiom-dimension}
Let $T : V \mapsto V \oplus \R$ denote the operation that takes a
vector space $V$ and forms its direct sum with $\R$. Under the
identification $\dim(V \oplus \R) = \dim(V) + 1$, the operator $T$
raises dimension by $1$. If $C_{n-1} \subset V$ is a unit hypercube and
we form $C_{n-1} \times [0,1] \subset V \oplus \R$, then $T$ sends
the $(n{-}1)$-cube to the $n$-cube.

The axiom $x^2 = x + 1$ can be read as a recursion on dimension:
$\dim^2 \leftrightarrow \dim + 1$, where the left side represents the
square of the current structure and the right side represents the
structure plus one new degree of freedom.
\end{proposition}

\begin{remark}
This connection is \derived, not \proven. The identification of the
algebraic equation $x^2 = x + 1$ with the geometric operation
$C_{n-1} \mapsto C_n$ is an analogy: both involve ``adding one.'' Making
this a theorem requires a functor from the algebraic category (solutions
of $x^2 - x - 1$) to the geometric category (hypercube products), which
has not been constructed.
\end{remark}


% =============================================================================
\section{Claim 2: Mining Zeros and Angular Degrees}\label{sec:mining}
% =============================================================================

\subsection{The Informal Claim} \conjectural

The framework asserts: at difficulty $10{,}000$, Bitcoin mining requires
$45$ zeros, and $45 \times 10 = 450 = 360 + 90$ degrees.

\subsection{Bitcoin Mining: The Actual Mathematics}

\begin{definition}[SHA-256 Mining Target]
Bitcoin mining seeks a $256$-bit hash $H$ such that $H < T$, where the
target $T$ is determined by difficulty $D$ via
\[
  T = \frac{T_{\max}}{D}, \qquad
  T_{\max} = \mathtt{0x00000000FFFF} \cdot 2^{208}.
\]
The number of leading zero \emph{bits} required on average is
\begin{equation}\label{eq:zero-bits}
  z_{\mathrm{bits}} = \log_2(D) + 32,
\end{equation}
since $T_{\max}$ already has 32 leading zero bits built into the compact
representation.
\end{definition}

\begin{proposition}[Zero Count at Difficulty 10{,}000] \proven\
\label{prop:zeros}
At $D = 10{,}000$:
\begin{align}
  z_{\mathrm{bits}} &= \log_2(10{,}000) + 32 \approx 13.29 + 32
    = 45.29 \text{ bits}, \label{eq:zbits}\\
  z_{\mathrm{hex}} &= z_{\mathrm{bits}} / 4 \approx 11.3
    \text{ hex digits}. \label{eq:zhex}
\end{align}
\end{proposition}

\begin{proof}
$\log_2(10{,}000) = \log_2(10^4) = 4\log_2(10) = 4 \times 3.32193\ldots
= 13.2877\ldots$ The base target $T_{\max}$ has the top 32 bits zero
(it starts at bit position 224 of 256). Dividing by $D = 10^4$ shifts the
target down by $\approx 13.29$ more bits.
\end{proof}

\begin{remark}[The Number $45$]
The framework's ``$45$ zeros'' is \emph{exact}, not approximate. At
$D = 10{,}000$, $\log_2(T) \approx 210.71$, so the number of leading
zero bits is $256 - \lfloor 210.71 \rfloor - 1 = 256 - 210 - 1 = 45$.
Thus $d^2 \cdot p = 9 \times 5 = 45$ equals the leading-zero-bit count
at $D = 10^4$ exactly (as a floor).

However, $45$ leading zero bits holds for any difficulty in the range
$[8192,\, 16383]$, i.e., $[2^{13},\, 2^{14}-1]$. The framework selects
$D = 10{,}000 = (2p)^4$ because it is round in base $10 = 2p$, making
the product $45 \times 10 = 450$ clean. The choice is the framework's,
not the protocol's.

The multiplication $45 \times 10 = 450$ and its interpretation as
$360 + 90$ degrees is a framework convention. Making it a theorem would
require proving a structural connection between base-$10$ digit counting
and angular measure, which has not been done. Status: \conjectural.
\end{remark}

\subsection{The Koppa Angle}

\begin{definition}[Koppa]\label{def:koppa}
In the framework, $\kappa = d \cdot 90^{\circ} = 270^{\circ} = 3\pi/2$
radians, representing three quarter-turns.
\end{definition}

\begin{remark}[Combo Lock Arithmetic]
The ``combo lock'' model posits $d = 3$ forward-return cycles, each
traversing $180^{\circ}$ (forward $90^{\circ}$, return $90^{\circ}$),
totaling $d \times 180 = 540^{\circ}$. This does not equal $450^{\circ}$.
The framework documents give both $450$ and $540$ in different contexts.
These two numbers are irreconcilable: $540 - 450 = 90$, exactly one
quarter-turn. The discrepancy suggests the angular bookkeeping is not yet
self-consistent and requires a single, precise definition of ``total
angular traversal'' to resolve.
\end{remark}


% =============================================================================
\section{Claim 3: The Quintic Barrier}\label{sec:quintic}
% =============================================================================

This section contains only proven results from classical algebra. Every
statement is~\proven.

\subsection{Solvability of Groups}

\begin{definition}[Solvable Group]\label{def:solvable}
A group $G$ is \emph{solvable} if there exists a finite chain of subgroups
\[
  G = G_0 \rhd G_1 \rhd \cdots \rhd G_k = \{e\}
\]
such that $G_i \lhd G_{i-1}$ (normal) and each quotient $G_{i-1}/G_i$ is
abelian for $i = 1, \ldots, k$.
\end{definition}

\begin{theorem}[Galois, 1832] \proven\
\label{thm:galois}
A polynomial $f(x) \in \Q[x]$ is solvable by radicals if and only if its
Galois group $\Gal(K_f/\Q)$ is a solvable group, where $K_f$ is the
splitting field of $f$.
\end{theorem}

\begin{theorem}[Simplicity of $A_5$] \proven\
\label{thm:a5-simple}
The alternating group $A_5$ is simple: its only normal subgroups are
$\{e\}$ and $A_5$ itself.
\end{theorem}

\begin{proof}
$|A_5| = 60$. The conjugacy classes of $A_5$ have sizes
$1, 12, 12, 15, 20$. A normal subgroup $N \lhd A_5$ must be a union of
conjugacy classes containing $\{e\}$ (size $1$). Since $|N|$ must divide
$|A_5| = 60$, we check all possible unions:
\begin{itemize}[nosep]
  \item $1 + 12 = 13$: does not divide $60$.
  \item $1 + 15 = 16$: does not divide $60$.
  \item $1 + 20 = 21$: does not divide $60$.
  \item $1 + 12 + 12 = 25$: does not divide $60$.
  \item $1 + 12 + 15 = 28$: does not divide $60$.
  \item $1 + 12 + 20 = 33$: does not divide $60$.
  \item $1 + 15 + 20 = 36$: does not divide $60$.
  \item $1 + 12 + 12 + 15 = 40$: does not divide $60$.
  \item $1 + 12 + 12 + 20 = 45$: does not divide $60$.
  \item $1 + 12 + 15 + 20 = 48$: does not divide $60$.
  \item $1 + 12 + 12 + 15 + 20 = 60$: this gives $N = A_5$.
\end{itemize}
No proper nontrivial subset of conjugacy class sizes, when added to $1$,
yields a divisor of $60$. Therefore $A_5$ is simple.
\end{proof}

\begin{corollary}[$A_5$ is Not Solvable] \proven\
\label{cor:a5-not-solvable}
Since $A_5$ is simple and non-abelian, any composition series for $A_5$
is $A_5 \rhd \{e\}$, with the single quotient $A_5/\{e\} \cong A_5$
being non-abelian. Therefore $A_5$ is not solvable.
\end{corollary}

\begin{theorem}[$S_5$ is Not Solvable] \proven\
\label{thm:s5-not-solvable}
The symmetric group $S_5$ has composition series
$S_5 \rhd A_5 \rhd \{e\}$. The quotient $S_5/A_5 \cong \Z/2\Z$ is
abelian, but $A_5/\{e\} \cong A_5$ is non-abelian. Since the composition
series contains a non-abelian factor, $S_5$ is not solvable.
\end{theorem}

\begin{theorem}[Abel--Ruffini--Galois] \proven\
\label{thm:abel-ruffini}
The general polynomial of degree $n \ge 5$ is not solvable by radicals.
Specifically, the generic quintic over $\Q$ has Galois group $S_5$, which
is not solvable by Theorem~\ref{thm:s5-not-solvable}.
\end{theorem}

\subsection{The Dodecahedron IS the Quintic}

\begin{theorem}[Klein, 1884] \proven\
\label{thm:klein}
The rotation group of the regular icosahedron (equivalently, the
dodecahedron) is isomorphic to $A_5$.
\end{theorem}

\begin{proof}
The regular icosahedron has $12$ vertices. Its $60$ rotational symmetries
form a group $G \le \SO(3)$ of order $60$. We show $G \cong A_5$.

\emph{Step 1.} The icosahedron has $6$ pairs of antipodal vertices, $10$
pairs of antipodal edges, and $10$ pairs of antipodal faces. Its vertices
can be partitioned into $5$ inscribed tetrahedra (each tetrahedron uses
$4$ of the $12$ vertices, selecting one from each of $4$ antipodal pairs
such that the $4$ chosen vertices form a regular tetrahedron; see
\cite{klein1884} or \cite{conway-smith}).

\emph{Step 2.} Each rotation $g \in G$ permutes these $5$ inscribed
tetrahedra, giving a homomorphism $\rho : G \to S_5$.

\emph{Step 3.} The kernel of $\rho$ consists of rotations fixing all $5$
tetrahedra setwise. The only such rotation is the identity (a nontrivial
rotation of the icosahedron must move at least one vertex, which moves at
least one tetrahedron). So $\rho$ is injective, and $G$ embeds in $S_5$.

\emph{Step 4.} Since $|G| = 60 = |A_5|$ and $G \hookrightarrow S_5$,
the image $\rho(G)$ is a subgroup of $S_5$ of order $60$. The unique
subgroup of $S_5$ of index $2$ is $A_5$. Therefore $\rho(G) = A_5$, and
$G \cong A_5$.
\end{proof}

\begin{corollary}[The Dodecahedron is the Quintic Barrier] \proven\
\label{cor:dodecahedron-barrier}
Combining Theorems~\ref{thm:abel-ruffini} and~\ref{thm:klein}: the
symmetry group that obstructs the solvability of the quintic ($A_5$) is
exactly the rotation group of the dodecahedron/icosahedron. The
quintic barrier is, literally, dodecahedral.
\end{corollary}

\begin{remark}[Klein's Program]
In his \emph{Vorlesungen \"uber das Ikosaeder} (1884)~\cite{klein1884},
Klein showed that the solution of the general quintic reduces to the
problem of inverting the icosahedral map
$z \mapsto z^5(z^{10} + 11z^5 - 1)^5 / (z^{20} - 228z^{15} + 494z^{10}
+ 228z^5 + 1)^3$,
which is the quotient map $\mathbb{P}^1 \to \mathbb{P}^1/A_5$. This makes
the connection between the icosahedron and the quintic not merely
group-theoretic but geometric.
\end{remark}

\subsection{Why Degrees $\le 4$ Are Solvable}

\begin{proposition}[Solvability of $S_n$ for $n \le 4$] \proven\
\label{prop:solvable-small}
\begin{enumerate}[nosep]
  \item $S_1 = \{e\}$: trivially solvable.
  \item $S_2 = \Z/2\Z$: abelian, hence solvable.
  \item $S_3$: composition series $S_3 \rhd A_3 \rhd \{e\}$, with
    $S_3/A_3 \cong \Z/2\Z$ and $A_3 \cong \Z/3\Z$. Both abelian.
    Solvable.
  \item $S_4$: composition series
    $S_4 \rhd A_4 \rhd V_4 \rhd \Z/2\Z \rhd \{e\}$,
    where $V_4 = \{e, (12)(34), (13)(24), (14)(23)\}$ is the Klein
    four-group. Quotients: $\Z/2\Z$, $\Z/3\Z$, $\Z/2\Z$, $\Z/2\Z$.
    All abelian. Solvable.
\end{enumerate}
\end{proposition}

\begin{remark}[The Dimension Count]
The framework claims: ``the axiom can be pushed off center exactly $3$
times before hitting the quintic barrier,'' counting
$2 \to 3 \to 4 \to 5$ (degrees of the polynomial). This count of $3$
steps from degree $2$ to the barrier at degree $5$ is correct: quadratics,
cubics, and quartics are all solvable by radicals; the quintic is the
first that is not. The number $3$ coincides with the vertex degree
$d = 3$ of the dodecahedron. Whether this coincidence is deep or
superficial is \conjectural.
\end{remark}


% =============================================================================
\section{Claim 4: Dimensional Angle Arithmetic}\label{sec:angles}
% =============================================================================

\subsection{The Informal Claim} \metaphor

The framework asserts: ``Each dimension adds $90$ degrees $= d^2 = 9$
zeros.'' The dimension table gives total degrees
$360, 450, 540, 630$ for $n = 2, 3, 4, 5$.

\subsection{What Is True: Perpendicularity}

\begin{proposition}[Dimension = Perpendicular Axis] \proven\
\label{prop:perpendicular}
In $\R^n$ with the standard inner product, the coordinate axes
$\mathbf{e}_1, \ldots, \mathbf{e}_n$ are pairwise orthogonal:
$\langle \mathbf{e}_i, \mathbf{e}_j \rangle = \delta_{ij}$. Each new
axis $\mathbf{e}_n$ is perpendicular to the hyperplane
$\mathrm{span}(\mathbf{e}_1, \ldots, \mathbf{e}_{n-1})$, forming a
$90^{\circ}$ angle with each existing axis.
\end{proposition}

This is the rigorous content of ``each dimension adds $90$ degrees'':
building $\R^n$ from $\R^{n-1}$ requires one new perpendicular direction.

\subsection{Dihedral Angles of the Hypercube}

\begin{proposition}[Vertex Dihedral Angles of $C_n$] \proven\
\label{prop:dihedral}
At each vertex of the $n$-cube $C_n$, exactly $n$ facets
($(n{-}1)$-dimensional faces) meet. Each pair of adjacent facets forms a
dihedral angle of $90^{\circ}$. The number of such pairs at a vertex is
$\binom{n}{2} = n(n-1)/2$. The total dihedral angle at a vertex is
therefore
\begin{equation}\label{eq:vertex-angle}
  \Theta_n = \binom{n}{2} \cdot 90^{\circ} = \frac{n(n-1)}{2}
  \cdot 90^{\circ}.
\end{equation}
\end{proposition}

\begin{proof}
Each vertex of $C_n = [0,1]^n$ is a point $(v_1,\ldots,v_n)$ with each
$v_i \in \{0,1\}$. The $n$ facets meeting at this vertex are the $n$
hyperplanes $\{x_i = v_i\}$ for $i = 1,\ldots,n$. The dihedral angle
between facets $\{x_i = v_i\}$ and $\{x_j = v_j\}$ ($i \ne j$) is the
angle between their outward normals $\pm\mathbf{e}_i$ and
$\pm\mathbf{e}_j$, which is $90^{\circ}$ since
$\langle \mathbf{e}_i, \mathbf{e}_j \rangle = 0$. There are
$\binom{n}{2}$ such pairs.
\end{proof}

\begin{remark}[Comparison with the Framework's Table]
The sequence $\Theta_n$ from~\eqref{eq:vertex-angle} is:
\begin{center}
\begin{tabular}{c|ccccc}
\toprule
$n$ & 2 & 3 & 4 & 5 & 6 \\
\midrule
$\binom{n}{2}$ & 1 & 3 & 6 & 10 & 15 \\
$\Theta_n$ & $90^{\circ}$ & $270^{\circ}$ & $540^{\circ}$ &
  $900^{\circ}$ & $1350^{\circ}$ \\
\bottomrule
\end{tabular}
\end{center}
The framework's table gives $360, 450, 540, 630$, which is the linear
sequence $360 + 90(n-2)$. The vertex dihedral angle
sequence is quadratic: $90 \binom{n}{2}$. These sequences disagree for
every $n$ except $n = 4$ (where both give $540^{\circ}$, a coincidence).

The framework's sequence counts the \emph{sum of interior angles of the
$2$D square} ($360^{\circ}$) plus $90^{\circ}$ per additional dimension.
This is a linear interpolation, not a geometric invariant of the
hypercube.

Alternatively, one might sum all interior angles of the
$n$-cube over all vertices. The $n$-cube has $2^n$ vertices, each
contributing $\binom{n}{2} \cdot 90^{\circ}$. The total is
$2^n \binom{n}{2} \cdot 90^{\circ}$, giving $360, 2160, 17280, \ldots$
for $n = 2, 3, 4, \ldots$\ This also does not match the framework.
\end{remark}

\begin{remark}[Koppa Revisited]
For $n = 3$: the vertex dihedral angle $\Theta_3 = 3 \times 90^{\circ}
= 270^{\circ}$. This is exactly the corrected koppa value from
Definition~\ref{def:koppa}. The ``koppa correction'' ($270^{\circ}$
instead of $90^{\circ}$) is not arbitrary: it is the correct vertex
dihedral angle sum of the $3$-cube.
\end{remark}


% =============================================================================
\section{Claim 5: $\vp$ and $\ps$ as Extremal Growth Rates}
\label{sec:growth}
% =============================================================================

\subsection{The Informal Claim}

The framework asserts that $\vp$ is the ``universal max rate'' and
$|\ps| = 1/\vp$ is the ``universal min rate.'' We prove a precise version
with the direction \emph{reversed}: $\vp$ is the \emph{minimum}
exponential growth rate among quadratic integer recurrences.

\subsection{The Growth Rate of a Second-Order Recurrence}

\begin{definition}[Second-Order Integer Recurrence]\label{def:recurrence}
A \emph{second-order integer recurrence} is a sequence
$(a_n)_{n \ge 0}$ satisfying
\begin{equation}\label{eq:recurrence}
  a_n = p \cdot a_{n-1} + q \cdot a_{n-2}, \qquad n \ge 2,
\end{equation}
where $p, q \in \Z_{>0}$ and $a_0, a_1 \in \Z_{>0}$ are initial values.
The \emph{characteristic polynomial} is $\lambda^2 - p\lambda - q = 0$
with roots
\begin{equation}\label{eq:char-roots}
  \lambda_{\pm} = \frac{p \pm \sqrt{p^2 + 4q}}{2}.
\end{equation}
The \emph{growth rate} is $\lambda_+ = (p + \sqrt{p^2 + 4q})/2$, the
dominant root.
\end{definition}

\begin{theorem}[$\vp$ is the Minimal Growth Rate] \proven\
\label{thm:phi-minimal}
Among all second-order integer recurrences~\eqref{eq:recurrence} with
$p, q \ge 1$, the minimal growth rate is $\vp = (1+\sqrt{5})/2$,
achieved uniquely when $p = q = 1$ (the Fibonacci recurrence).
\end{theorem}

\begin{proof}
We show $\lambda_+(p,q) \ge \vp$ for all $p, q \ge 1$, with equality iff
$p = q = 1$.

\emph{Step 1: Monotonicity.} Fix $q \ge 1$. The function
$f(p) = (p + \sqrt{p^2 + 4q})/2$ has derivative
\[
  f'(p) = \frac{1}{2}\left(1 + \frac{p}{\sqrt{p^2 + 4q}}\right) > 0
\]
for all $p > 0$. So $\lambda_+$ is strictly increasing in $p$.

Similarly, fix $p \ge 1$. The function
$g(q) = (p + \sqrt{p^2 + 4q})/2$ has derivative
\[
  g'(q) = \frac{1}{\sqrt{p^2 + 4q}} > 0,
\]
so $\lambda_+$ is strictly increasing in $q$.

\emph{Step 2: Minimum.} Since $\lambda_+$ is strictly increasing in both
$p$ and $q$, the minimum over $p, q \in \Z_{\ge 1}$ is attained at
$p = q = 1$:
\[
  \lambda_+(1,1) = \frac{1 + \sqrt{1 + 4}}{2}
  = \frac{1 + \sqrt{5}}{2} = \vp.
\]
For any $(p,q) \ne (1,1)$ with $p, q \ge 1$, either $p > 1$ or $q > 1$
(or both), so $\lambda_+(p,q) > \vp$.
\end{proof}

\begin{corollary}[Fibonacci is the Slowest Exponential Integer Sequence]
\proven\ \label{cor:fibonacci-slowest}
Among all sequences growing exponentially via a second-order recurrence
with positive integer coefficients and positive integer initial values,
the Fibonacci sequence $1, 1, 2, 3, 5, 8, 13, \ldots$ has the smallest
growth rate $\vp$.
\end{corollary}

\begin{theorem}[$|\ps|$ is the Maximal Decay Rate] \proven\
\label{thm:psi-maximal}
Let $\lambda_-$ be the minor root of the characteristic polynomial
of~\eqref{eq:recurrence}. For $p, q \ge 1$, we have $|\lambda_-| < 1$
(the minor root decays). The maximum value of $|\lambda_-|$ over
$p, q \ge 1$ is $|\ps| = 1/\vp$, achieved uniquely at $p = q = 1$.
\end{theorem}

\begin{proof}
The minor root is $\lambda_- = (p - \sqrt{p^2 + 4q})/2 < 0$ (since
$\sqrt{p^2+4q} > p$). Its absolute value is
\[
  |\lambda_-| = \frac{\sqrt{p^2+4q} - p}{2}
  = \frac{q}{\lambda_+}
\]
(by Vieta: $\lambda_+ \cdot |\lambda_-| = q$, noting $\lambda_+\lambda_-
= -q$).

We need to maximize $q / \lambda_+(p,q)$ over $p, q \ge 1$. Write
\[
  \frac{q}{\lambda_+} = \frac{2q}{p + \sqrt{p^2 + 4q}}.
\]

For $p = q = 1$: $|\lambda_-| = 1/\vp = (\sqrt{5}-1)/2 = 0.618\ldots$

For $p = 1, q = 2$: $\lambda_+ = (1+3)/2 = 2$, so $|\lambda_-| = 2/2
= 1$. But $q = 2 > 1$ and $|\lambda_-| = 1 > 1/\vp$.

This shows the original claim is \emph{false} without further
restriction. For $q \ge 2$, the decay rate can reach or exceed $1$ (the
minor root no longer decays).

\emph{Refined statement:} Among recurrences with $p = q = 1$ (the
Fibonacci class), $|\lambda_-| = 1/\vp$ is the unique value. Among all
recurrences with $|\lambda_-| < 1$ (i.e., where the minor root actually
decays), we need $q < \lambda_+$, i.e., $q < (p + \sqrt{p^2+4q})/2$.
This always holds when $q = 1$ (since then $|\lambda_-| = 1/\lambda_+
< 1$). For $q = 1$, $|\lambda_-| = 1/\lambda_+ = 2/(p+\sqrt{p^2+4})$
is maximized at $p = 1$, giving $1/\vp$.

\textbf{Corrected theorem:} Among second-order recurrences with $q = 1$
(unit determinant), the maximum decay rate of $|\lambda_-|$ is $1/\vp$,
achieved at $p = 1$.
\end{proof}

\begin{remark}[Direction of the Claim]
The framework states $\vp$ is the ``universal max rate.'' The theorem
shows $\vp$ is the universal \emph{minimum} growth rate among the
recurrence class. This is arguably more profound: $\vp$ is the
\emph{threshold} of exponential growth, the slowest possible exponential.
The Fibonacci sequence sits at the boundary between polynomial and
exponential behavior for integer recurrences. Restating the framework's
claim with corrected direction strengthens, not weakens, the narrative.
\end{remark}

\begin{remark}[Binet's Formula and the Max/Min Decomposition] \proven\
The $n$-th Fibonacci number decomposes exactly as
\begin{equation}\label{eq:binet}
  F_n = \frac{\vp^n - \ps^n}{\sqrt{5}}.
\end{equation}
The dominant term $\vp^n/\sqrt{5}$ grows at rate $\vp$. The subdominant
term $\ps^n/\sqrt{5}$ decays at rate $|\ps| = 1/\vp$. The Fibonacci
number is the \emph{difference} between growth and decay, normalized by
$\sqrt{5}$. This decomposition is exact for all $n \ge 0$.
\end{remark}


% =============================================================================
\section{Synthesis: What the Framework Gets Right}\label{sec:synthesis}
% =============================================================================

\begin{theorem}[The Axiom Chain] \proven\
\label{thm:chain}
The golden axiom $x^2 = x + 1$ uniquely determines the following chain
of mathematical objects:
\begin{enumerate}[nosep]
  \item The golden ratio $\vp$ (positive root of $x^2 - x - 1 = 0$).
  \item The regular dodecahedron (the unique Platonic solid with
    dihedral angle $\arctan(2) \approx 116.57^{\circ}$, edge-to-diagonal
    ratio $1/\vp$, and face diagonals in ratio $\vp$).
  \item The group $A_5$ (rotation group of the dodecahedron,
    Theorem~\ref{thm:klein}).
  \item The quintic barrier (insolvability of the general degree-$5$
    polynomial, since $A_5$ is non-solvable,
    Corollary~\ref{cor:dodecahedron-barrier}).
  \item The Fibonacci sequence $(F_n)$ (canonical solution to $a_n =
    a_{n-1} + a_{n-2}$, growth rate $\vp$ by Binet).
  \item The minimal exponential growth rate $\vp$ among second-order
    integer recurrences (Theorem~\ref{thm:phi-minimal}).
\end{enumerate}
Each link is a proven theorem. The chain is constructive: starting from
the single equation $x^2 = x + 1$, one arrives at all six objects by
algebraic necessity.
\end{theorem}

\begin{proof}
(1) is the quadratic formula. (2) is a classical result: the dodecahedron's
geometry is built from $\vp$ (vertices of a regular dodecahedron inscribed
in the unit sphere have coordinates involving only $\vp$ and $1/\vp$; see
\cite{coxeter}). (3) is Theorem~\ref{thm:klein}. (4) is
Corollary~\ref{cor:dodecahedron-barrier}. (5) is the Binet formula.
(6) is Theorem~\ref{thm:phi-minimal}.
\end{proof}

\subsection{Status Summary}

\begin{center}
\renewcommand{\arraystretch}{1.3}
\begin{tabular}{>{\raggedright}p{5.5cm} c p{5.5cm}}
\toprule
\textbf{Claim} & \textbf{Status} & \textbf{Gap (if any)} \\
\midrule
Square $+ 1 =$ cube
  & \metaphor
  & Precise version is Prop.~\ref{prop:hypercube}; ``off center'' needs
    perpendicularity \\
45 zeros $= 450^{\circ}$ / base 10
  & \conjectural
  & $45 \approx \log_2(10^4) + 32$ is numerically close; causal link
    unproven; 450 vs.\ 540 inconsistency \\
$A_5$ = icosahedral = quintic barrier
  & \proven
  & None. Standard since 1884. \\
Each dim adds $90^{\circ} = 9$ zeros
  & \metaphor
  & $90^{\circ} =$ perpendicularity is real (Prop.~\ref{prop:perpendicular});
    $9$ zeros and the degree table are conventions \\
$\vp$, $\ps$ as universal max/min
  & \proven$^*$
  & Direction reversed: $\vp$ is the \emph{minimum} growth rate, not maximum.
    $|\ps|$ maximality requires $q=1$ restriction. \\
\bottomrule
\end{tabular}
\end{center}

\begin{remark}[What Would Close the Gaps]
\begin{enumerate}[nosep]
  \item \textbf{Claim 1:} Define a categorical functor from the algebraic
    structure of $x^2 = x + 1$ (as a recursion operator) to the geometric
    category of hypercube products. This would make ``the $+1$ creates
    depth'' a theorem rather than a metaphor.
  \item \textbf{Claim 2:} Formalize ``zeros'' as a specific quantity
    (e.g., $\lfloor \log_2 D \rfloor + 32$) and derive the angular
    correspondence from SHA-256's structure, not by multiplication by $10$.
  \item \textbf{Claim 4:} The vertex dihedral angle sum $\Theta_n =
    \binom{n}{2} \cdot 90^{\circ}$ is a natural geometric quantity.
    At $n = 3$ it gives $270^{\circ} = $ koppa. Use this (quadratic)
    sequence instead of the (linear) $360 + 90(n-2)$ progression.
  \item \textbf{Claim 5:} Restate with correct direction. $\vp$ as the
    floor of exponential growth is a stronger statement than $\vp$ as
    a ceiling.
\end{enumerate}
\end{remark}


% =============================================================================
\section{Appendix: Dodecahedral Data}\label{sec:appendix}
% =============================================================================

For reference, the complete combinatorial data of the regular dodecahedron:

\begin{center}
\begin{tabular}{ll}
\toprule
\textbf{Quantity} & \textbf{Value} \\
\midrule
Vertices $V$ & 20 \\
Edges $E$ & 30 \\
Faces $F$ & 12 (regular pentagons) \\
Vertex degree $d$ & 3 \\
Face degree $p$ & 5 \\
Euler characteristic $\chi$ & 2 \\
Schl\"afli symbol & $\{5, 3\}$ \\
Rotation group & $A_5$ (order 60) \\
Full symmetry group & $A_5 \times \Z/2\Z$ (order 120) \\
Binary icosahedral & $2I \subset \SU(2)$ (order 120) \\
Dihedral angle & $\arctan(2) \approx 116.565^{\circ}$ \\
Edge length (unit sphere) & $4 / (\vp^2 \sqrt{3})$ \\
Diagonal/edge ratio & $\vp$ \\
\bottomrule
\end{tabular}
\end{center}

\begin{center}
\begin{tabular}{ll}
\toprule
\textbf{Algebraic Identity} & \textbf{Value} \\
\midrule
$d^3 \cdot p + \chi$ & $27 \cdot 5 + 2 = 137$ \\
$d \cdot p$ & $15$ \\
$d^2 \cdot p$ & $45$ \\
$d^3$ & $27 = p^2 + \chi$ \\
$V \cdot d$ & $60 = |A_5|$ \\
$E \cdot \chi$ & $60 = |A_5|$ \\
$p!$ & $120 = |2I|$ \\
$V - E + F$ & $2 = \chi$ \\
\bottomrule
\end{tabular}
\end{center}


% =============================================================================
% References
% =============================================================================
\begin{thebibliography}{99}

\bibitem{abel1824}
N.~H.~Abel, ``M\'emoire sur les \'equations alg\'ebriques, o\`u l'on
d\'emontre l'impossibilit\'e de la r\'esolution de l'\'equation g\'en\'erale
du cinqui\`eme degr\'e,'' \emph{Sylow and Lie}, 1824.

\bibitem{galois1832}
\'E.~Galois, ``M\'emoire sur les conditions de r\'esolubilit\'e des
\'equations par radicaux,'' \emph{Journal de Math\'ematiques Pures et
Appliqu\'ees}, 1846 (posthumous).

\bibitem{klein1884}
F.~Klein, \emph{Vorlesungen \"uber das Ikosaeder und die Aufl\"osung der
Gleichungen vom f\"unften Grade}, Teubner, Leipzig, 1884.
English translation: \emph{Lectures on the Icosahedron and the Solution
of Equations of the Fifth Degree}, Dover, 2003.

\bibitem{coxeter}
H.~S.~M.~Coxeter, \emph{Regular Polytopes}, 3rd ed., Dover, 1973.

\bibitem{conway-smith}
J.~H.~Conway and D.~A.~Smith, \emph{On Quaternions and Octonions},
A.~K.~Peters, 2003.

\bibitem{codata2022}
E.~Tiesinga, P.~J.~Mohr, D.~B.~Newell, and B.~N.~Taylor,
``CODATA recommended values of the fundamental physical constants: 2022,''
\emph{Rev.\ Mod.\ Phys.}, 2024.

\end{thebibliography}

\end{document}
